% The composite operators of FermiWick and their definitions (included by the showcase and by the manual).
\begin{center}\small
\renewcommand{\arraystretch}{1.25}
\begin{tabular}{lll}\toprule
FermiWick & operator & definition \\ \midrule
\texttt{nOp[i, s]} & $n_{i\sigma}$ & $c^\dagger_{i\sigma}c_{i\sigma}$ \\
\texttt{nTotOp[i]}, \texttt{nOp[i]} & $n_i$ & $n_{i\uparrow}+n_{i\downarrow}$ \\
\texttt{SzOp[i]} & $S^z_i$ & $\tfrac12\,(n_{i\uparrow}-n_{i\downarrow})$ \\
\texttt{SpOp[i]}, \texttt{SmOp[i]} & $S^+_i$, $S^-_i$ & $c^\dagger_{i\uparrow}c_{i\downarrow}$, \ $c^\dagger_{i\downarrow}c_{i\uparrow}$ \\
\texttt{SpinOperator[i, a]} & $S^a_i$ & $\tfrac12\sum_{\sigma\sigma'}c^\dagger_{i\sigma}\,\sigma^a_{\sigma\sigma'}\,c_{i\sigma'}$, \ $a=1,2,3$ \\
\texttt{DOp[i]} & $D_i$ (double occupancy) & $n_{i\uparrow}n_{i\downarrow}$ \\
\texttt{Bdag[i]}, \texttt{B[i]} & $B^\dagger_i$, $B_i$ (doublon) & $c^\dagger_{i\uparrow}c^\dagger_{i\downarrow}$, \ $c_{i\downarrow}c_{i\uparrow}$ \\
\texttt{EmptyProjector[i]} & $P^0_i$ & $(1-n_{i\uparrow})(1-n_{i\downarrow})=1-n_i+D_i$ \\
\texttt{SingleOccupationProjector[i]} & $P^1_i$ & $n_i-2D_i$ \\
\texttt{GutzwillerProjector[i]} & $P^G_i$ & $\sqrt2\,[\cos\theta\,P^1_i+\sin\theta\,(P^0_i+D_i)]$ \\
\texttt{SdotS[i, j]} & $\mathbf S_i\cdot\mathbf S_j$ & $S^z_iS^z_j+\tfrac12(S^+_iS^-_j+S^-_iS^+_j)$ \\
\texttt{SdotScrossS[i, j, k]} & $\mathbf S_i\cdot(\mathbf S_j\times\mathbf S_k)$ & $\sum_{abc}\epsilon_{abc}\,S^a_iS^b_jS^c_k$ \\ \bottomrule
\end{tabular}
\end{center}
Products of these on one site reduce with the algebra of the fermions (e.g.\ $D_in_i=2D_i$, $(S^z_i)^2=\tfrac14n_i-\tfrac12D_i$);
operators on several sites and products with symbolic sites stay as products. \texttt{GutzwillerProjector} is in
\texttt{FermiWick`Gutzwiller`}; one-site operators on several sites are products, \texttt{DOp[i, j]} $=D_iD_j$.
